\begin{table}[H]
\raggedright
\caption{Extended six-node model: autoregressive effects and effects on pain.}
\label{tab:sup-extended}
\scriptsize
\begingroup
\setlength{\tabcolsep}{4pt}
\renewcommand{\arraystretch}{1.08}
\noindent
\begin{tabular*}{\textwidth}{@{\extracolsep{\fill}}lcccc@{}}
\toprule
Effect (lag-1) & Estimate & SE & $p$ & Random SD \\
\midrule
Threat value $\rightarrow$ Pain intensity & 0.096 & 0.027 & <.001 & 0.137 \\
Attention to pain $\rightarrow$ Pain intensity & 0.035 & 0.019 & .064 & 0.069 \\
Goal engagement $\rightarrow$ Pain intensity & 0.008 & 0.013 & .544 & 0.037 \\
Activity valence $\rightarrow$ Pain intensity & -0.005 & 0.014 & .736 & 0.041 \\
Goal efficacy $\rightarrow$ Pain intensity & -0.003 & 0.013 & .817 & 0.0 \\
Pain intensity $\rightarrow$ Pain intensity & 0.346 & 0.025 & <.001 & 0.153 \\
Threat value $\rightarrow$ Threat value & 0.298 & 0.028 & <.001 & 0.171 \\
Attention to pain $\rightarrow$ Attention to pain & 0.179 & 0.026 & <.001 & 0.156 \\
Goal engagement $\rightarrow$ Goal engagement & 0.138 & 0.024 & <.001 & 0.128 \\
Goal efficacy $\rightarrow$ Goal efficacy & 0.07 & 0.018 & <.001 & 0.027 \\
Activity valence $\rightarrow$ Activity valence & 0.065 & 0.019 & .001 & 0.021 \\
\bottomrule
\end{tabular*}
\par\vspace{3pt}\noindent\parbox{\textwidth}{\scriptsize\textit{Note.} The model adds two further goal-directed activity characteristics, momentary efficacy and valence. The core four-node edges were essentially unchanged relative to the primary benchmark (see the sensitivity-agreement table).}
\endgroup
\end{table}
